% =============================================================================
% Kapitel 3: Methodik und KI-gestützte Entwicklung
% =============================================================================

\chapter{Methodik und KI-gestützte Entwicklung}
\label{chap:methodik}

Große Sprachmodelle unterstützten die Auswertung des RFC-Primärtexts, den Entwurf und die Implementierung,
die Analyse der Messkampagnen und die Prüfung von Zwischenergebnissen. Wissenschaftlich belastbar ist dieses
Vorgehen nur, wenn die Arbeit offenlegt, wie Modellausgaben geprüft werden und wer die Ergebnisse verantwortet.
Diese Regeln müssen vor der Ergebnisbewertung feststehen.

\section{Anlass des KI-Einsatzes und veröffentlichte Fälle}
\label{sec:externe-faelle}

Diese Arbeit entsteht als Einzelarbeit und verbindet die Auswertung einer neuen Norm, systemnahe
Programmierung, formale Modellierung und Messkampagnen. Die Arbeit legt den Einsatz von KI-Werkzeugen offen
und beschreibt die Regeln zur Prüfung ihrer Ergebnisse. Zur Einordnung dienen neun
veröffentlichte Fälle, in denen KI-Systeme, darunter Sprachmodelle, an Forschungs- und
Entwicklungsergebnissen beteiligt waren. Sie bilden keine systematische Literaturübersicht, belegen keine
einheitliche Prüfkette und erlauben keine Aussage über Häufigkeit oder Wirksamkeit eines Verfahrens. Die
folgende Auflistung nennt nur die in den Quellen ausdrücklich dokumentierten Rollen: den Beitrag des
KI-Systems, die prüfende Instanz und die Grenze der jeweiligen Prüfung.

\begin{itemize}
  \item \textbf{AlphaFold:} Ein Vorhersagemodell bestimmte Proteinstrukturen; der blinde CASP14-Vergleich
    mit experimentell ermittelten Strukturen prüfte die Ergebnisse, eine allgemeine Prüfkette entsteht
    daraus nicht \cite{jumper2021alphafold}.
  \item \textbf{Davies u.\,a.:} Ein KI-System schlug Hypothesen und Kandidatenmuster für mathematische
    Probleme vor; beteiligte Mathematiker prüften durch Gegenbeispielsuche und führten den klassischen
    Beweis \cite{davies2021intuition}.
  \item \textbf{FunSearch:} Ein Sprachmodell erzeugte ausführbare Programmkandidaten. Ein Prüfprogramm
    bewertete jeden Kandidaten; das Ergebnis gilt nur für die im Programm festgelegte Bewertungsfunktion
    \cite{romera2024funsearch}.
  \item \textbf{AlphaProof:} Ein Modell erzeugte formale Beweiskandidaten, die der Kern des
    Beweisassistenten Lean einzeln nachprüfte; jeder Beweis gilt für die formalisierte Aussage und ihre
    Annahmen \cite{alphaproof2026}.
  \item \textbf{Knuths Zyklen:} Ein Modell konstruierte die Lösung eines Zyklenproblems; Knuth führte den
    Beweis anschließend selbst, und unabhängig davon lag bereits eine veröffentlichte Lean-Formalisierung
    von Morrison vor \cite{knuth2026cycles}.
  \item \textbf{Zeta-Schranke:} Ein Modell lieferte ein mathematisches Ergebnis und formalisierte es
    gemeinsam mit dem Anthropic-Mitarbeiter Eric Easley in Lean. Zwei Anthropic-Mathematiker und zwei
    externe Fachleute prüften das Ergebnis; ein Blindvergleich ist nicht dokumentiert \cite{anthropic2026zeta}.
  \item \textbf{Einheitsabstand:} Ein Modell lieferte einen Gegenbeweis zu einer Vermutung der diskreten
    Geometrie; externe Mathematiker prüften ihn und veröffentlichten eine verdichtete, von Menschen
    verifizierte Fassung \cite{openai2026unitdistance}.
  \item \textbf{C-Parameter:} Ein Modell steuerte unter Aufsicht Rechnungen, Numerik und Manuskriptteile
    bei; Physiker prüften mit einer Gegenprobe durch Monte-Carlo-Simulationen, und der Preprint legt den
    Einsatz offen \cite{schwartz2026cparameter}.
  \item \textbf{Bun-Migration:} Bei der Migration von Bun von Zig nach Rust erzeugten und prüften Modelle
    derselben Familie den Quelltext. Trotz getrennter Rollen blieb ein Fehler unentdeckt. Nach dem Merge
    wurde undefiniertes Verhalten gemeldet und das Prüfwerkzeug Miri in die kontinuierliche Integration
    aufgenommen \cite{bun-in-rust,bun-pr30412,bun-issue30719}.
\end{itemize}

Die Fälle belegen nur, dass die Prüfinstanz zum Artefakt passen muss. Normtext, ausführbare Tests, formale
Beweise, Messdaten und menschliche Freigabe haben verschiedene Aufgaben; ihre tatsächliche Nutzung muss
dokumentiert sein. Diese Arbeit folgt außerdem der ausdrücklichen Offenlegung des C-Parameter-Falls
(\secref{sec:ki-offenlegung}) und versteht auch eine mehrstufige Prüfung nicht als Garantie. Sie wendet
denselben Maßstab auf sich selbst an: Jede Artefaktart erhält eine passende Prüfinstanz, und die
Arbeitsteilung zwischen Autor und Werkzeugen steht fest, bevor Ergebnisse bewertet werden.

\section{Forschungsdesign und KI-gestützte Arbeitsweise}
\label{sec:forschungsdesign}
\label{sec:rfc-auswertung}

Die Untersuchung verbindet die Auswertung der Norm, die Referenzimplementierung und eine empirische
Pfaduntersuchung. \figref{fig:forschungsdesign} zeigt, wie die beiden Forschungsfragen auf diesen Teilen
aufbauen:

\begin{figure}[H]
  \centering
  \begin{tikzpicture}[
      x=1cm, y=1cm,
      >={Stealth[length=5pt, width=4pt]},
      block/.style={draw, semithick, rounded corners=3pt, align=center,
        font=\scriptsize\linespread{1.0}\selectfont, inner sep=4pt},
      fluss/.style={->, draw=sqdunkel, thin},
      flusstext/.style={font=\scriptsize, inner sep=1.5pt, fill=white},
    ]
    % Rahmen und Ebenen
    \draw[draw=sqgrau, semithick, rounded corners=6pt] (0.15,0.85) rectangle (14.85,-5.75);
    \draw[draw=sqgrau, dashed, thin] (0.15,-1.15) -- (14.85,-1.15);
    \draw[draw=sqgrau, dashed, thin] (0.15,-3.3) -- (14.85,-3.3);
    \node[font=\scriptsize, text=sqdunkel] at (0.95,0)     {Norm};
    \node[font=\scriptsize, text=sqdunkel] at (0.95,-2.2)  {Software};
    \node[font=\scriptsize, text=sqdunkel] at (0.95,-4.62) {Messung};
    % Norm-Ebene
    \node[block, draw=sqgoldrand, fill=sqgold, minimum height=0.8cm] (quellen) at (3.1,0) {Normquellen};
    \draw[draw=sqblaurand, semithick, fill=sqblau, rounded corners=2pt] (5.55,0.55) rectangle (8.85,-0.55);
    \node[font=\scriptsize\bfseries] at (7.2,0.24) {Anforderungsmodell};
    \foreach \i/\f in {0/white, 1/sqblau, 2/white, 3/sqblau}
      \draw[draw=sqblaurand, thin, fill=\f] (5.8+\i*0.7,-0.38) rectangle (6.5+\i*0.7,-0.02);
    \node[block, draw=sqroserand, fill=sqrose, minimum height=0.8cm] (ffeins) at (12.8,0)
      {\textbf{FF1}\\ \textbf{technische Erfüllbarkeit}};
    \draw[fluss] (quellen) -- node[flusstext, above] {ableiten} (5.55,0);
    \draw[fluss] (8.85,0) -- node[flusstext, above] {Maßstab} (ffeins);
    % Software-Ebene
    \draw[draw=sqtealrand, semithick, fill=white] (5.7,-1.75) rectangle (8.7,-2.65);
    \draw[draw=sqtealrand, semithick, fill=white] (5.55,-1.98) rectangle (5.85,-2.16);
    \draw[draw=sqtealrand, semithick, fill=white] (5.55,-2.24) rectangle (5.85,-2.42);
    \node[font=\scriptsize] at (7.25,-2.2) {Rust-Bibliothek};
    \draw[fluss] (7.2,-0.55) -- node[flusstext, anchor=west, xshift=2pt] {umsetzen} (7.2,-1.75);
    \draw[fluss] (8.7,-2.2) -| node[flusstext, pos=0.25, above] {Prüfobjekt} (ffeins.south);
    % Messungs-Ebene
    \draw[fluss] (7.2,-2.65) -- node[flusstext, anchor=west, xshift=2pt] {Messmittel} (7.2,-3.9);
    \draw[draw=sqblaurand, semithick, fill=white, rounded corners=2pt] (3.6,-3.9) rectangle (11.2,-5.35);
    \fill[sqblau] (3.64,-3.94) rectangle (11.16,-4.32);
    \node[font=\scriptsize\bfseries] at (7.4,-4.13) {Netzpfad};
    \node[block, draw=sqdunkel, fill=sqblau, minimum height=0.55cm, inner sep=3pt]
      (spcap) at (4.85,-4.85) {Sender-PCAP};
    \node[block, draw=sqblaurand, fill=white, minimum height=0.55cm, inner sep=3pt]
      (epcap) at (9.9,-4.85) {Empfänger-PCAP};
    \foreach \i in {0,1,2}
      \draw[draw=sqdunkel, thin, fill=white] (6.75+\i*0.75,-4.85) circle (0.18);
    \draw[fluss] (spcap) -- (6.55,-4.85);
    \draw[fluss] (6.95,-4.85) -- (7.3,-4.85);
    \draw[fluss] (7.7,-4.85) -- (8.05,-4.85);
    \draw[fluss] (8.45,-4.85) -- (epcap);
    \node[block, draw=sqroserand, fill=sqrose, minimum height=0.8cm] (ffzwei) at (12.8,-4.62)
      {\textbf{FF2}\\ \textbf{Durchlässigkeit}};
    \draw[fluss] (11.2,-4.62) -- (ffzwei);
  \end{tikzpicture}
  \caption{Aufbau der Untersuchung und Zuordnung der Forschungsfragen}
  \label{fig:forschungsdesign}
\end{figure}

Aus \figref{fig:forschungsdesign} geht die Doppelrolle der Referenzimplementierung hervor: Für FF1 ist sie der
Prüfgegenstand, und das aus den Normquellen abgeleitete Anforderungsmodell bildet den Maßstab; für FF2 dient
dieselbe Implementierung als Messmittel. Beide Bewertungen bleiben getrennt: Eine konforme Implementierung
belegt keine Zustellbarkeit entlang realer Netzpfade, und eine unveränderte Zustellung belegt keine
RFC-Konformität.

Die Forschungsfragen benötigen deshalb verschiedene Untersuchungseinheiten und Entscheidungsregeln.
\tabref{tab:operationalisierung} stellt sie gegenüber:

\begin{table}[htbp]
  \centering
  \caption{Operationalisierung der Forschungsfragen}
  \label{tab:operationalisierung}
  \small
  \begin{tabularx}{\textwidth}{@{}L{0.7cm}L{3.2cm}YL{4.4cm}@{}}
    \toprule
    \textbf{FF} & \textbf{Untersuchungseinheit} & \textbf{Nachweis} & \textbf{Bewertung und Grenze} \\
    \midrule
    FF1
      & jede anwendbare Zeile des Arbeitsindex aus \secref{sec:anforderungsmatrix} im Projektumfang
        (eine normative Aussage oder ein Bündel davon)
      & je nach Status: Normfundstelle sowie Quelltext und Prüfung oder Beleg einer technischen Grenze
      & vollständig, teilweise oder nicht erfüllbar; nicht anwendbar bleibt unbewertet \\
    FF2
      & kontrolliertes Szenario je Pfad, Richtung und Messzeitpunkt
      & Paketmitschnitte an Sender und Empfänger, Sendemanifest und dokumentierte maschinelle Auswertung
      & Surplus erhalten, verändert oder entfernt; Datagramm verworfen oder Lauf ungültig. Gilt nur für Pfad,
        Richtung und Messzeitpunkt \\
    \bottomrule
  \end{tabularx}
\end{table}

Aus \tabref{tab:operationalisierung} geht hervor, dass FF1 je anwendbarer Anforderung und FF2 je
kontrolliertem Szenario entschieden wird. Ein Szenario kann je ein Kontroll- und ein Optionsdatagramm,
mehrere Wiederholungen oder bei FRAG mehrere übertragene Datagramme enthalten.
Für FF1 wird zunächst die Anwendbarkeit bestimmt. Anwendbare \texttt{MUST}- und
\texttt{SHOULD}-Aussagen gehen in die Bewertung ein; eine Abweichung von \texttt{SHOULD} benötigt eine
dokumentierte Begründung. Eine \texttt{MAY}-Funktion wird nur bewertet, wenn sie für die
Referenzimplementierung gewählt wurde.

Wie diese Regeln auf die einzelnen Anforderungen angewandt werden, legt \secref{sec:ff1-kriterien} fest; die
Auswertungsregeln der Messläufe für FF2 stehen in \secref{sec:testkonzept}. Neben der Erfüllbarkeit wird der
Umsetzungsstand berichtet. Dafür werden die Bibliothek und die Software, die sie aufruft, gemeinsam
bewertet; die zugehörige Regel definiert \secref{sec:soll-ist}.

Der Maßstab für FF1 entsteht aus dem Primärtext. Normativ ist RFC~9868; RFC~768 liefert ergänzend das
UDP-Format, RFC~1071 als informative Referenz die Rechenregeln der Prüfsummen. Die Schlüsselwörter
\texttt{MUST}, \texttt{SHOULD} und \texttt{MAY} werden nach BCP~14\footnote{BCP steht für Best Current
Practice. BCP~14 umfasst RFC~2119 und dessen Aktualisierung RFC~8174.} gelesen. Nur die großgeschriebenen
Formen tragen die dort festgelegte besondere Bedeutung; normative Aussagen können die Schlüsselwörter aber
auch auslassen \cite{rfc2119,rfc8174}. RFC~9868 stellt Absätzen mit solchen Wörtern zusätzlich das
Zeichenpaar \texttt{>{}>} voran. Die eigene Zählung von 103 markierten Absätzen ist deshalb nur eine
Kontrollzahl und keine vollständige Anforderungsliste. Für die Belegform gilt: Normbezüge nennen den
maßgeblichen RFC-Abschnitt im Fließtext oder im Quellenverweis; bei BCP-14-Aussagen sowie bei Zahlen und
Grenzwerten trägt der Quellenverweis die Abschnittsangabe immer.

Jede so gewonnene Aussage wird nach demselben Muster erfasst: Fundstelle, Normstufe, Anwendbarkeit und
Projektumfang. Dieses Muster ist zugleich das Zeilenschema des Anforderungskatalogs im
Repository der Implementierung, sodass jede Katalogzeile ihren RFC-Abschnitt mitführt. RFC~9869 und die Errata
zu RFC~9868 werden gesondert behandelt. RFC~9869 beschreibt die Pfad-MTU-Bestimmung (DPLPMTUD, \textit{Datagram
Packetization Layer Path MTU Discovery}) mit UDP-Optionen und liegt außerhalb des Implementierungsumfangs
\cite{rfc9869}. Von den Errata ist das verifizierte technische EID~8834 angewendet; es löst den Widerspruch
zwischen den Abschnitten~10 und 14 zur Behandlung überlanger Optionen \cite{rfc9868-eid8834}. Das
redaktionelle EID~8708 bleibt für eine spätere Ausgabe des RFC zurückgestellt \cite{rfc9868-eid8708}.
\figref{fig:rfc-auswertung} fasst diesen Weg zusammen:

\begin{figure}[H]
  \centering
  \begin{tikzpicture}[
      x=1cm, y=1cm,
      >={Stealth[length=5pt, width=4pt]},
      block/.style={draw, semithick, rounded corners=3pt, align=center,
        font=\scriptsize\linespread{1.0}\selectfont, inner sep=4pt},
      fluss/.style={->, draw=sqdunkel, thin},
      flusstext/.style={font=\scriptsize, inner sep=1.5pt, fill=white, align=center},
    ]
    % Hauptfluss
    \node[block, draw=sqgoldrand, fill=sqgold] (quellen) at (1.75,0)
      {\textbf{Normquellen}\\ RFC 9868 (normativ)\\ RFC 768 (ergänzend)\\ RFC 1071 (informativ)};
    \node[block, draw=sqblaurand, fill=sqblau] (lesung) at (5.6,0)
      {\textbf{BCP-14-Lesung}\\ \texttt{>{}>}-markierte Absätze\\ Kontrollzahl: 103\\ zwei Zählverfahren};
    \node[block, draw=sqblaurand, fill=sqblau] (erfassung) at (9.45,0)
      {\textbf{Erfassung je Aussage}\\ Fundstelle, Normstufe,\\ Anwendbarkeit,\\ Projektumfang};
    \node[block, draw=sqroserand, fill=sqrose] (matrix) at (13.3,0)
      {\textbf{Abgleich mit der}\\ \textbf{Konformitäts-}\\ \textbf{matrix}\\ Maßstab für FF1};
    \draw[fluss] (quellen) -- (lesung);
    \draw[fluss] (lesung) -- (erfassung);
    \draw[fluss] (erfassung) -- (matrix);
    % Errata-Detail
    \node[block, draw=sqgrau, fill=sqhell] (errata) at (1.75,1.75)
      {\textbf{Errata-Stand}\\ EID 8834 angewendet,\\ EID 8708 zurückgestellt};
    \draw[draw=sqgrau, dashed, thin] (errata.south) -- (quellen.north);
    % Abgrenzung und Kontrolle
    \node[block, draw=sqgrau, fill=sqhell, dashed] (ausserhalb) at (1.75,-2.35)
      {RFC 9869 (DPLPMTUD):\\ außerhalb des Umfangs};
    \draw[draw=sqgrau, dashed, thin] (quellen.south) --
      node[flusstext] {getrennt gehalten} (ausserhalb.north);
    \node[block, draw=sqgrau, fill=sqhell] (erstzaehlung) at (5.6,-2.35)
      {Erstzählung: 96 Absätze,\\ zweites Verfahren\\ deckte den Fehler auf};
    \draw[draw=sqgrau, dashed, thin] (lesung.south) -- (erstzaehlung.north);
    % Zweiter Zufluss
    \node[block, draw=sqgoldrand, fill=sqgold] (vorgaben) at (10.1,-2.35)
      {Projekt-, Schnittstellen- und\\ Plattformvorgaben};
    \draw[fluss] (vorgaben.east) -| (matrix.south);
    \node[flusstext, anchor=east, align=right] at (13.18,-1.55) {ausdrücklich\\ gekennzeichnet};
  \end{tikzpicture}
  \caption{Von den Normquellen zum geprüften Anforderungsbestand}
  \label{fig:rfc-auswertung}
\end{figure}

\figref{fig:rfc-auswertung} zeigt die beiden Grundlagen des Maßstabs: die Anforderungen des RFC unter
Berücksichtigung der Errata und die davon getrennt erfassten eigenen Projekt-, Schnittstellen- und
Plattformvorgaben. Die eigenen Vorgaben sind ausdrücklich gekennzeichnet. Diese Trennung verhindert, dass
eine lokale Entwurfsentscheidung nachträglich als Forderung des RFC erscheint. Die Zahl der markierten
Absätze wurde mit zwei Zählverfahren geprüft: Die Erstzählung ergab 96 Absätze; erst das zweite Verfahren
deckte den Fehler auf. Zählung und Einordnung erfolgten beim Lesen des RFC-Texts. Ein Werkzeug zur
automatischen Zerlegung des Texts wurde dafür nicht eingesetzt.

Diese Erfassung liegt nicht im Quelltext, sondern in versionierten Textdokumenten im
Repository der Implementierung. Sie sind ausdrücklich an die Werkzeuge gerichtet: Die zentrale Datei mit
Arbeitsvorgaben
nennt in ihrem ersten Satz das Modellwerkzeug und menschliche Mitwirkende als
Adressaten.\footnote{\raggedright Repository der Implementierung, Stand vom 14.~August 2026: \path{CLAUDE.md},
\mbox{\path{docs/requirements.md}}, \path{docs/plan/ROADMAP.md} sowie die Schrittdateien unter
\path{docs/plan/steps/}.} Die Dokumente halten den Wissensstand über die Arbeitssitzungen hinweg fest und
dienen als Maßstab für die Prüfung der Modellausgaben. Ein Abgleich am Primärtext fand vier falsche
RFC-Angaben in vier Dokumenten mit Arbeitsvorgaben zugleich.

Die Versionsgeschichte schreibt diese Vorgaben dem Autor zu: Die beiden ersten Commits des Repositorys
legten Roadmap, Schrittdateien, Anforderungskatalog sowie Architektur- und Formatreferenz an. Die Roadmap
beschreibt einen Gesamtplan mit menschlicher Freigabe: Jeder Schritt soll mit einem geprüften Commit
abgeschlossen werden. Der Autor legt damit Plan, Maßstab und Abschlusskriterien fest und gibt jedes
Ergebnis frei; die Werkzeuge arbeiten innerhalb dieser Vorgaben. Aus den Sammelcommits lässt sich die
Reihenfolge der Arbeiten innerhalb eines Schritts nicht rekonstruieren. Die Commit-Historie zeigt außerdem
nicht, welchen Beitrag Autor und Werkzeug zu einer einzelnen Änderung geleistet haben.

Jeder Hauptschritt bekommt eine eigene Datei nach festem Schema mit Ziel,
Anforderungen, Lean-Pflichten, Plan, Aufgaben und Abschlusskriterien; alle Schrittdateien lagen vor dem
ersten Funktionsschritt vor. Wo Lean-Pflichten bestehen, verlangen neun Schrittdateien wörtlich die
Spezifikation vor der Umsetzung und den Beweis danach; \figref{fig:stepzyklus} zeigt den Kernzyklus:

\begin{figure}[H]
  \centering
  \begin{tikzpicture}[
      x=1cm, y=1cm,
      >={Stealth[length=5pt, width=4pt]},
      block/.style={draw, semithick, rounded corners=3pt, align=center, text width=3.0cm,
        font=\scriptsize\linespread{1.0}\selectfont, inner sep=4pt},
      fluss/.style={->, draw=sqdunkel, thin},
      ausnahme/.style={->, draw=sqgrau, dashed, thin},
      flusstext/.style={font=\scriptsize, inner sep=1.5pt, fill=white, align=center},
    ]
    \node[block, draw=sqgoldrand, fill=sqgold] (schritt) at (1.75,0)
      {\textbf{Schritt wählen}\\ Schrittdatei: Plan und\\ Abschlusskriterien};
    \node[block, draw=sqgrau, fill=sqhell] (lean) at (5.6,0)
      {bei Lean-Pflicht:\\ Spezifikation vor dem\\ Code, Beweise danach};
    \node[block, draw=sqblaurand, fill=sqblau] (umsetzung) at (9.45,0)
      {\textbf{Umsetzung}\\ Code und Tests zusammen;\\ neue oder geänderte\\ Paketverarbeitung:\\
       Property-Tests und\\ Fuzz-Ziele ergänzen};
    \node[block, draw=sqblaurand, fill=sqblau] (gate) at (13.3,0)
      {\textbf{lokale Prüfung}\\ \texttt{pre-pr.sh}\\ 18 Prüfläufe, Abbruch\\ beim ersten Fehler};
    \node[block, draw=sqblaurand, fill=sqblau] (pr) at (13.3,-2.7)
      {\textbf{Pull Request}\\ CI: Format, Lint,\\ Tests, Lean};
    \node[block, draw=sqgrau, fill=sqhell] (ki) at (9.45,-2.7)
      {zwei angeforderte\\ KI-Prüfungen: Qualität;\\ RFC-Semantik};
    \node[block, draw=sqroserand, fill=sqrose] (freigabe) at (5.6,-2.7)
      {\textbf{menschliche Freigabe}\\ ein Sammelcommit\\ auf \texttt{main}};
    \draw[fluss] (schritt) -- (lean);
    \draw[fluss] (lean) -- (umsetzung);
    \draw[fluss] (umsetzung) -- (gate);
    \draw[fluss] (gate) -- (pr);
    \draw[fluss] (pr) -- (ki);
    \draw[ausnahme] (ki) -- (freigabe);
    \draw[ausnahme] (ki.north) -- node[flusstext, midway, right] {Befund beheben} (umsetzung.south);
    \draw[fluss] (freigabe.north) -- ++(0,0.55) -| (schritt.south);
    \node[flusstext] at (3.68,-1.5) {nächster Schritt};
  \end{tikzpicture}
  \caption[Kernzyklus eines Implementierungsschritts]{Kernzyklus eines Implementierungsschritts bis
    Schritt 18 (Ausnahmen gestrichelt)}
  \label{fig:stepzyklus}
\end{figure}

Sprachmodelle sind in den Untersuchungen Werkzeuge, aber keine bewertende Instanz. Maßgeblich sind die
Normquellen, die Messdaten, ausführbare Prüfungen und die Freigabe durch den Autor. \tabref{tab:rollen}
trennt die Rollen von Autor und Werkzeugen nach Aufgabe und Grenze:

\begin{table}[htbp]
  \centering
  \caption{Rollen bei Erzeugung, Prüfung und Freigabe}
  \label{tab:rollen}
  \small
  \begin{tabularx}{\textwidth}{@{}L{3.0cm}L{4.4cm}Y@{}}
    \toprule
    \textbf{Beteiligter} & \textbf{Aufgabe} & \textbf{Grenze} \\
    \midrule
    Autor
      & legt Plan, Anforderungskatalog, Architektur- und Formatvorgaben sowie die Schrittdateien fest;
        entscheidet, gewichtet, gibt frei und verantwortet
      & jedes Ergebnis gilt erst nach seiner Freigabe \\
    Claude Code (Anthropic)
      & setzt innerhalb dieser Vorgaben um: Entwürfe, Quelltext, Tests und Analysen
      & Modellausgaben sind keine Nachweise \\
    Codex (OpenAI)
      & getrennter Zweitprüfer
      & Herstellertrennung beweist keine Unabhängigkeit \\
    \bottomrule
  \end{tabularx}
\end{table}

Beide Modellwerkzeuge verbinden ein Sprachmodell mit Dateizugriff, Kommandoausführung und Repository. Der
Einsatz von Modellen verschiedener Hersteller soll das Risiko verringern, dass ihre Fehlermuster vollständig
übereinstimmen. Unabhängige Prüfungen sind damit jedoch nicht belegt. Eine Prüfung durch dasselbe Modell ist
keine getrennte Zweitprüfung. Wenn der Zweitprüfer selbst Änderungen vornimmt, braucht er dafür einen
eigenen, vom Autor freigegebenen Auftrag. Diese Änderungen zählen nicht als getrennte Zweitprüfung des
bearbeiteten Ergebnisses.

Die Art der Prüfung richtet sich nach der Art der Aussage. Protokollaussagen werden gegen RFC-Primärtext und
Errata geprüft. Aussagen über die Implementierung werden durch Quelltext und ausführbare Prüfungen gestützt.
Pfadbefunde benötigen Paketmitschnitte und Kampagnenprotokolle. Literaturbehauptungen werden an der jeweiligen
Primärquelle geprüft. Ein Hashwert belegt nur die Unverändertheit eines Artefakts, nicht die Wahrheit seines
Inhalts.

Die Arbeit mit den Modellen folgt drei schriftlich festgehaltenen Regeln:

\begin{enumerate}
  \item Modellergebnisse werden gegen die jeweilige Primärquelle geprüft, nie gegen eigene Zusammenfassungen.
  \item Aufträge werden neutral formuliert, ohne das Modell in eine gewünschte Richtung zu lenken.
  \item Widerspruch wird ausdrücklich ermutigt: Ein Modell, das \zitat{nein} sagt oder Unsicherheit meldet,
    ist wertvoller als eines, das gefällig zustimmt.
\end{enumerate}

Diese Regeln beschreiben den Sollprozess; welche Prüfungen tatsächlich stattfanden, wird nicht aus ihm
abgeleitet, sondern aus den Protokollen rekonstruiert (\secref{sec:reproduzierbarkeit}). Warum diese Trennung
nötig ist, zeigt der folgende Abschnitt.

\section{Bekannte Schwächen der Sprachmodelle}
\label{sec:gefahrenmodell}

\tabref{tab:sprachmodell-schwaechen} ordnet den bekannten Schwächen der Sprachmodelle die beobachtete
Wirkung und die methodische Antwort zu.

\begin{table}[H]
  \centering
  \caption{Schwächen der Sprachmodelle und methodische Antworten}
  \label{tab:sprachmodell-schwaechen}
  \small
  \begin{tabularx}{\textwidth}{@{}L{3.0cm}L{5.5cm}Y@{}}
    \toprule
    \textbf{Schwäche} & \textbf{Wirkung und Beleg} & \textbf{Antwort der Methodik} \\
    \midrule
    Halluzination
      & überzeugend formulierter, aber falscher Inhalt; Trainings- und Bewertungsverfahren können Raten
        gegenüber eingestandener Unsicherheit belohnen \cite{kalai2025hallucinate}
      & Prüfung an der Primärquelle, nicht an eigenen Zusammenfassungen \\
    Gefälligkeit (Sycophancy)
      & erkennbare Erwartungen können Widerspruch unterdrücken; das Verhalten tritt bereits bei reinen
        Vortrainingsmodellen auf, steigt bei größeren Modellen, bleibt nach Training mit menschlicher
        Rückmeldung bestehen und verursachte in Versuchen über elf Modelle messbare Schäden
        \cite{perez2022evals,cheng2026sycophancy}
      & neutrale Aufträge, ausdrücklich erwünschter Widerspruch und getrennter Zweitprüfer \\
    Verfügbarkeit
      & Cloud-Dienste können ausfallen, gedrosselt oder verändert werden; wechselnde Modelle und
        nichtdeterministische Antworten verhindern die vollständige Reproduktion einer Modellausgabe
      & dauerhafte, ohne den Dienst prüfbare Artefakte: Quelltext, Tests, Protokolle und Messdaten \\
    \bottomrule
  \end{tabularx}
\end{table}

Halluzination und Gefälligkeit können zusammenwirken, wenn eine falsche Aussage unwidersprochen bleibt.
Die Gegenmaßnahmen setzen deshalb an Auftrag, Prüfinstanz und Artefakten an. Keine Aussage gilt, weil ein
Modell sie behauptet. Konformitätsmatrix und Testarchitektur prüfen die Implementierung; auch sie belegen
nur die untersuchten Eigenschaften und keinen allgemeinen Korrektheitsbeweis.

Menschen und Modelle müssen die Verfahrensregeln aus der Tabelle anwenden. Ausführbare Prüfungen und der
Beweisprüfer liefern nach ihrem Start automatisch ein Ergebnis.
Wie diese Prüfungen für FF1 aufgebaut sind, legt der folgende Abschnitt fest.

\section{Testarchitektur für FF1}
\label{sec:pruefkette}

Die aktuelle Prüfkette wird lokal gestartet, verbindet aber lokale und entfernte Prüfungen. Sie ist der
Sollzustand vor dem Öffnen oder Aktualisieren eines Pull Requests. Historisch entstand diese Kette
schrittweise.

Dieser Abschnitt beschreibt, wie die Implementierung während der Entwicklung für FF1 geprüft wird.
Ein bestandener Entwicklungstest sagt nichts über reale Netzpfade aus, und eine Pfadmessung ersetzt keinen
Entwicklungstest. Die Pfaduntersuchung für FF2 verwendet die fertigen Programme, Sendemanifeste sowie
Paketmitschnitte an Sender und Empfänger; ihre Ergebnisklassen legt bereits
\tabref{tab:operationalisierung} fest, das Pfad- und Auswertungsmodell folgt in \secref{sec:ff2-pfadmodell},
die Belegregeln und die Auswertung der Messläufe in \secref{sec:testkonzept}. Welche Auswerter dabei
eingesetzt werden und was sie prüfen, beschreibt \secref{subsec:mitschnittwerkzeuge}; die Bewertungsregeln
beider Untersuchungen bleiben getrennt.

\subsection{Prüfmittel und ihre Reichweite}
\label{subsec:pruefmittel}

Zu den festen Prüfungen gehören Formatierung, statische Analyse, Dokumentationsbau sowie Unit- und
Integrationstests. Ein Property-Test formuliert eine
Eigenschaft, die für alle Eingaben einer Klasse gelten soll, und prüft sie an vielen maschinell erzeugten
Fällen statt an ausgewählten Beispielen \cite{proptest}.

Ein Fuzz-Ziel ist eine Funktion, die eine
beliebige Bytefolge entgegennimmt und damit die echten Verarbeitungspfade der Bibliothek aufruft, vom
Zerlegen über das Serialisieren bis zum Aufbau ganzer Datagramme. Der abdeckungsgeleitete Fuzzer erzeugt
und verändert solche Eingaben fortlaufend und bevorzugt jene, die neue Programmpfade erreichen
\cite{cargo-fuzz,libfuzzer}. Fuzzing kann Fehler finden, aber keine Fehlerfreiheit beweisen. Es läuft als
lokale Pflichtstufe und nicht in der kontinuierlichen Integration.

Lean ist ein Beweisassistent: Regeln werden als Modell formuliert, Aussagen darüber als Theoreme bewiesen, und ein
kleiner Kern prüft jeden Beweisschritt nach \cite{lean4}. Die Lean-Prüfung baut diese Modelle und
kontrolliert ihre Axiome. Bewiesen ist damit das handgeschriebene Modell mit seinen Annahmen, nicht der
Rust-Quelltext.

Eine besondere Gefahr sind Fehler, die Implementierung und Test gemeinsam haben. Sender und Empfänger der
Bibliothek verwenden getrennte Module für Serialisierung und Parsing, beruhen aber auf denselben
projektspezifischen Formatannahmen, etwa der gemeinsamen Tabelle der \texttt{Kind}-Werte und der
Längenkonstanten. Tests, die auf denselben fehlerhaften Annahmen beruhen, können den Entwurfsfehler deshalb
übersehen. Bei der Wire-Prüfung zeichnet \texttt{tcpdump} die tatsächlich gesendeten Bytes für eine
getrennte Auswertung auf. Ein getrenntes Python-Programm dekodiert IPv4, UDP und die UDP-Optionen und
berechnet die Prüfsummen.
\texttt{tshark} bestätigt zusätzlich die IPv4- und UDP-Felder. Die im Projekt verwendete Version 4.6.4
besitzt jedoch keinen Dissektor für UDP-Optionen.

\subsection{Prüfkette bis zur Freigabe}

Das lokale Prüfskript \texttt{pre-pr.sh} führt neun feste Prüfläufe aus: Formatierung, statische Analyse,
Dokumentationsbau, Tests mit Property-Fällen, Prüfung des eigenen Auswerters, Lean-Prüfung und drei
Prüfläufe für das Linux-Zielsystem. Diese umfassen Cross-Kompilierung und Tests, Raw-Socket-Integration mit
Root-Rechten und die Prüfung der tatsächlich gesendeten Pakete. Die Cross-Kompilierung erfolgt auf dem
Entwicklungsrechner; die Testprogramme laufen auf der lokalen virtuellen Maschine \texttt{achim}.
Für jedes der neun Fuzz-Ziele kommt ein weiterer Prüflauf hinzu.
Insgesamt ergeben sich die 18 Prüfläufe aus \figref{fig:stepzyklus}. Die kontinuierliche Integration (CI)
wiederholt Formatierung, statische Analyse und Unit-Tests. Nach deren Erfolg baut sie die Lean-Modelle und
stellt anschließend zwei getrennte Prüfanfragen: an Claude Code zu Codequalität und Sicherheit und an
Codex zur RFC-9868-Semantik.

Der aktuelle Sollprozess ist in \figref{fig:pruefprozess} als Sequenzdiagramm dargestellt: Jeder Beteiligte
erhält eine senkrechte Lebenslinie, die schmalen Balken zeigen seine aktiven Phasen, durchgezogene Pfeile sind
Aufträge und gestrichelte Pfeile Rückmeldungen.

\begin{figure}[H]
  \centering
  \begin{tikzpicture}[
      x=1cm, y=1cm,
      >={Stealth[length=5pt, width=4pt]},
      kopf/.style={draw, semithick, rounded corners=3pt, minimum height=0.62cm, inner xsep=5pt,
        font=\footnotesize},
      ruf/.style={->, draw=sqdunkel, thin},
      antwort/.style={->, draw=sqdunkel, thin, dashed},
      pfeiltext/.style={font=\scriptsize\linespread{1.0}\selectfont, inner sep=1.5pt, fill=white,
        anchor=south, align=center},
      notiz/.style={draw=sqgrau, fill=sqhell, rounded corners=2pt, thin, align=left,
        font=\scriptsize\linespread{1.0}\selectfont, inner sep=3pt},
    ]
    \def\xAutor{0.8}\def\xSkript{3.42}\def\xLinux{6.04}\def\xHub{8.66}\def\xCodex{11.28}\def\xClaude{13.9}
    % Lebenslinien
    \foreach \x in {\xAutor,\xSkript,\xLinux,\xHub,\xCodex,\xClaude}
      \draw[sqgrau, dashed, thin] (\x,-0.31) -- (\x,-9.7);
    % Köpfe
    \node[kopf, draw=sqgoldrand, fill=sqgold] at (\xAutor,0)  {Autor};
    \node[kopf, draw=sqblaurand, fill=sqblau] at (\xSkript,0) {Prüfskript};
    \node[kopf, draw=sqtealrand, fill=sqteal] at (\xLinux,0)  {Linux-Ziel};
    \node[kopf, draw=sqdunkel,   fill=sqhell] at (\xHub,0)    {GitHub};
    \node[kopf, draw=sqroserand, fill=sqrose] at (\xCodex,0)  {Codex};
    \node[kopf, draw=sqroserand, fill=sqrose] at (\xClaude,0) {Claude Code};
    % Aktivierungsbalken
    \draw[sqgoldrand, fill=sqgold, thin] (\xAutor-0.08,-1.05)  rectangle (\xAutor+0.08,-9.3);
    \draw[sqblaurand, fill=sqblau, thin] (\xSkript-0.08,-1.05) rectangle (\xSkript+0.08,-3.75);
    \draw[sqtealrand, fill=sqteal, thin] (\xLinux-0.08,-1.85)  rectangle (\xLinux+0.08,-3.05);
    \draw[sqdunkel,   fill=sqhell, thin] (\xHub-0.08,-4.3)     rectangle (\xHub+0.08,-9.3);
    \draw[sqroserand, fill=sqrose, thin] (\xCodex-0.08,-6.95)  rectangle (\xCodex+0.08,-7.55);
    \draw[sqroserand, fill=sqrose, thin] (\xClaude-0.08,-6.1)  rectangle (\xClaude+0.08,-8.1);
    % Lokale Prüfung
    \draw[ruf] (\xAutor+0.08,-1.05) -- node[pfeiltext] {\texttt{pre-pr.sh} starten} (\xSkript-0.08,-1.05);
    \draw[ruf] (\xSkript+0.08,-1.85) -- node[pfeiltext]
      {Tests ohne Root,\\ Raw-Socket-Test,\\ Wire-Prüfung} (\xLinux-0.08,-1.85);
    \node[notiz, anchor=east] (hoststufen) at (\xSkript-0.4,-2.4)
      {\textbf{Lokale Prüfungen}\\ Format, Analyse,\\ Tests, Property,\\ Fuzz, Lean,\\ Cross-Kompilierung};
    \draw[sqgrau, dashed, thin] (hoststufen.east) -- (\xSkript-0.08,-2.4);
    \draw[antwort] (\xLinux-0.08,-3.05) -- node[pfeiltext] {Ergebnisse} (\xSkript+0.08,-3.05);
    \draw[antwort] (\xSkript-0.08,-3.75) -- node[pfeiltext] {Gesamtergebnis} (\xAutor+0.08,-3.75);
    % Pull Request und CI
    \draw[ruf] (\xAutor+0.08,-4.3) -- node[pfeiltext] {Pull Request öffnen oder aktualisieren}
      (\xHub-0.08,-4.3);
    \node[notiz, anchor=west] (cistufen) at (\xHub+0.4,-4.95)
      {\textbf{CI-Stufen}\\ Rust-Prüfungen,\\ danach Lean-Bau};
    \draw[sqgrau, dashed, thin] (\xHub+0.08,-4.95) -- (cistufen.west);
    % Zweitprüfungen
    \draw[ruf] (\xHub+0.08,-6.1) -- node[pfeiltext] {Prüfanfrage Codequalität und Sicherheit}
      (\xClaude-0.08,-6.1);
    \draw[ruf] (\xHub+0.08,-6.95) -- node[pfeiltext, xshift=2.5pt] {Prüfanfrage\\ RFC-9868-Semantik}
      (\xCodex-0.08,-6.95);
    \draw[antwort] (\xCodex-0.08,-7.55) -- node[pfeiltext] {Befund} (\xHub+0.08,-7.55);
    \draw[antwort] (\xClaude-0.08,-8.1) -- node[pfeiltext] {Befund} (\xHub+0.08,-8.1);
    \draw[sqgrau, dashed, semithick, rounded corners=3pt] (0.1,-5.65) rectangle (14.8,-8.35);
    % Freigabe
    \draw[antwort] (\xHub-0.08,-8.7) -- node[pfeiltext] {Befunde} (\xAutor+0.08,-8.7);
    \draw[ruf] (\xAutor+0.08,-9.3) -- node[pfeiltext] {Freigabe und Merge} (\xHub-0.08,-9.3);
  \end{tikzpicture}
  \caption{Prüfprozess im Sollzustand}
  \label{fig:pruefprozess}
\end{figure}

Aus \figref{fig:pruefprozess} geht hervor, dass die Pfeilfolge nur die vorgesehene Reihenfolge vom Prüfskript
über den Pull Request und die Rust/Lean-CI bis zur Freigabe durch den Autor zeigt; eine technische
Merge-Bedingung ist sie nicht. Nach erfolgreicher Rust/Lean-CI werden getrennte Zweitprüfungen angefragt.

In dieser Arbeit gilt eine Prüfung nur dann als belegt, wenn sie protokolliert ist. Die Grenzen dieser
Protokollierung behandelt der folgende Abschnitt.

\section{Reproduzierbarkeit und Grenzen}
\label{sec:reproduzierbarkeit}

Aussagen über den Entstehungsprozess sind nur so belastbar wie ihre Protokollierung. Maßgeblich sind die
vorhandenen Arbeitsprotokolle, also die Repository-Historie, die Prüf- und Kampagnenprotokolle sowie die
Messdaten. Technische Ergebnisse sind reproduzierbar, soweit Quellstand, Werkzeuge, Eingaben und Ausgaben
archiviert sind. Die genaue Folge nichtdeterministischer Modellausgaben ist es nicht.

\chapref{chap:analyse} verwendet die hier festgelegten Quellen- und Bewertungsregeln, um aus RFC~9868 das
Anforderungsmodell für FF1 abzuleiten.
